% !TeX root = ../main.tex
\section{Conclusion}
\label{sec:conclusion}

A room-level graph formulation for conditional shear wall layout generation is presented, with rooms modeled as nodes and adjacencies as edges to better reflect architectural semantics and boundary-constrained placement. Conditioning via FiLM-style modulation enables density-regime control, and a dual-stream training procedure with counterfactual regularization is used to improve condition sensitivity when paired condition-specific supervision is limited.

Across 143 residential floor plans, the room-based representation improves prediction quality over edge-based baselines, achieving 0.597 Image IoU and yielding higher precision (+17.5 pp), which is consistent with reduced false positive wall placement. Ablation results indicate that explicit conditioning is the primary driver of controllability (CGS -0.148 when removed) and that the proposed training strategy provides a smaller but measurable benefit to conditional compliance. Case-based structural validation using finite element analysis further suggests that generated layouts can satisfy feasibility checks under the evaluated settings.

Future work should consider continuous condition formulations, representations that accommodate non-rectangular spaces without ad hoc preprocessing, and tighter integration of structural response metrics into training objectives for performance-aware generation. Extending the framework toward joint prediction of multiple boundary-defined structural elements is also a natural direction, but it requires systematic validation beyond the current scope.
